% =============================================================================
\chapter{Séance 6 — Maintenance opérationnelle et supervision}\label{chap:seance6}
% =============================================================================

\fiche{En tant que techniciens d'exploitation, contrôler l'état de la plateforme, fiabiliser la
maintenance, rendre les anomalies observables et conserver des preuves vérifiables : socle
obligatoire UC-01 à UC-06, compléments UC-C1 à UC-C4 et un défi avancé (UC-A3).}
{MIR Mahan, Jolan Allen, Kenzi Bendjelloul}{23/10/2026}{\fait{} — UC-05 \partiel}

\section{Déroulé et bilan}

La séance a commencé par un blocage : l'hôte Proxmox avait été redémarré et les VM, sans démarrage
automatique, étaient éteintes. Une fois les VM démarrées, un snapshot \code{s09_avant_seance6} a
été pris avant toute modification (figure~\ref{fig:s6-chrono}).

\illustration{01-chronologie.pdf}{1}{Déroulé réel de la séance 6 (en rouge : blocage et tests de panne contrôlés)}{fig:s6-chrono}

\begin{table}[H]
\centering
\caption{Bilan de la séance 6 par use case}\label{tab:s6-bilan}
\begin{tabularx}{\linewidth}{@{}P{3.6cm}LP{1.4cm}P{2.1cm}@{}}
\toprule
\textbf{Use case} & \textbf{Réalisation} & \textbf{Statut} & \textbf{Preuves} \\
\midrule
UC-01 État de santé & Fiche d'état initial datée avant toute modification ; 5 écarts relevés & \fait & \preuve{P-80}, \preuve{P-80c} \\
UC-02 Métriques & 4/4 cibles collectées ; écart de règles et vérification du playbook corrigés & \fait & \preuve{P-81} \\
UC-03 Tableau de bord & Dashboard v3 : unités, descriptions, données absentes distinguées de zéro & \fait & \preuve{P-85} \\
UC-04 Alerte & 2 alertes déclenchées puis résolues par des pannes contrôlées & \fait & \preuve{P-84} \\
UC-05 Incident & Arrêt brutal du DC : cause confirmée, correction préparée & \partiel & \preuve{P-86} \\
UC-06 Livrables & Code, CI, README, dossier ; validé / partiel / bloqué listé & \fait & \preuve{P-64} \\
UC-C1 Maintenance & Maintenance de LINUX01 entre deux contrôles de santé & \fait & \preuve{P-87} \\
UC-C2 Contrôle Ansible & \code{controle_sante.yml} en lecture seule & \fait & \preuve{P-82} \\
UC-C3 Tendance & Mesure ponctuelle comparée à 4 jours d'historique & \fait & \preuve{P-88} \\
UC-C4 Tâche planifiée & Timer systemd tracé dans le journal et dans Prometheus & \fait & \preuve{P-83} \\
UC-A3 Quotas (défi) & Quotas sur le disque de test \code{/data} & \fait & \preuve{P-89} \\
\bottomrule
\end{tabularx}
\end{table}

% -----------------------------------------------------------------------------
\section{UC-01 — État de santé vérifiable}

Le relevé initial (09:06) couvre l'identité, l'heure, les disques, les services, les unités en
échec, les erreurs récentes et un test fonctionnel par service (résolution DNS, page HTTP avec
contrôle du contenu, API Prometheus et Grafana, requête Active Directory). Sous Linux, il utilise
\code{hostnamectl}, \code{timedatectl}, \code{df}, \code{systemctl}, \code{journalctl},
\code{getent} et \code{curl} ; sous Windows, un script PowerShell en lecture seule exécuté par WinRM.

\begin{table}[H]
\centering
\caption{Écarts relevés par l'état initial et suite donnée}\label{tab:s6-ecarts}
\begin{tabularx}{\linewidth}{@{}P{0.4cm}P{4.6cm}LP{1.6cm}@{}}
\toprule
\textbf{\#} & \textbf{Écart} & \textbf{Cause et suite} & \textbf{Statut} \\
\midrule
1 & VM éteintes après le redémarrage de l'hôte & Pas de démarrage automatique ; VM démarrées à la main & Levé \\
2 & Jeton API Proxmox sans droit & Séparation des privilèges ; retirée par l'étudiant & Levé \\
3 & 4 règles d'alerte chargées au lieu de 5 & Copie du projet sur OPS01 non resynchronisée ; supervision redéployée & Corrigé \\
4 & \code{openipmi.service} en échec sur les VM Linux & Paquet recommandé installé par nos playbooks, sans matériel IPMI ; service désactivé par Ansible & Corrigé \\
5 & WIN-DC01 : événements 41 et 6008 & Arrêt non propre à chaque arrêt de l'hôte : incident UC-05 & Diagnostiqué \\
\bottomrule
\end{tabularx}
\end{table}

\illustration{P-80c-rendu-proxmox-snapshots-s09.png}{0.95}{Snapshots des trois VM, dont \code{s09_avant_seance6} (rendu de la sortie de l'API, \preuve{P-80c})}{fig:s6-snap}

% -----------------------------------------------------------------------------
\section{UC-02 — Métriques Linux et Windows collectables}

La collecte a été vérifiée de bout en bout : configuration chargée, connexions TCP établies vers
les exporters, endpoints \code{/metrics} en HTTP 200, puis valeurs lues dans la base de Prometheus
(figure~\ref{fig:s6-flux}).

\illustration{02-flux-supervision.pdf}{1}{Architecture de supervision après la séance 6 (en bleu : éléments ajoutés)}{fig:s6-flux}

\illustration{P-81-prometheus-cibles-up.png}{0.85}{Prometheus : 4/4 cibles \code{up} après correction (\preuve{P-81})}{fig:s6-cibles}

\attention{Défaut révélé : vérification des cibles qui réussissait sans cible}{%
Juste après un redémarrage de Prometheus, la liste des cibles est vide : l'assertion du playbook
réussissait alors sans rien vérifier (\og 0 cibles up \fg{}). Au passage suivant, une cible encore
\og unknown \fg{} la faisait échouer. Correction : attendre le premier scrape (8 tentatives de 5 s)
et exiger une liste non vide ; le cas défaillant a été rejoué avec succès (\preuve{P-81a} → \preuve{P-81b}).}

% -----------------------------------------------------------------------------
\section{UC-03 — Tableau de bord d'exploitation}

Le dashboard \og InfraReady — Serveurs \fg{} passe en version 3 : description d'usage, période de
6 heures, une description \og ce que le panneau permet de décider \fg{} sur chacun des 9 panneaux,
des unités explicites, et surtout la distinction entre une \textbf{donnée absente} (\og NON
COLLECTÉ \fg{}, courbes interrompues) et une \textbf{valeur réellement nulle}. Trois panneaux sont
ajoutés : contrôle planifié du site, maintenance Linux et évolution du stockage.

\illustration{P-85-grafana-dashboard-v3-exploitation.png}{0.82}{Dashboard Grafana v3 : le laboratoire était éteint avant 9 h, d'où les zones vides plutôt que des zéros (\preuve{P-85})}{fig:s6-dash}

% -----------------------------------------------------------------------------
\section{UC-04 et UC-C4 — Tâche planifiée et alertes testées}

L'incident de la phase H avait montré qu'un site indisponible ne déclenchait aucune alerte, puisque
seul l'exporter était surveillé. Un \textbf{timer systemd} lance désormais \code{check-nginx.sh}
chaque minute sur LINUX01 et publie le résultat pour Prometheus via le collecteur \og textfile \fg{}
de node\_exporter. Le service tourne en root \textbf{sans aucune capacité}, avec un système de
fichiers en lecture seule sauf le dossier des métriques.

\illustration{03-sequence-alerte.pdf}{0.95}{Chaîne de détection pendant le test 1 (nginx arrêté)}{fig:s6-seq}

\begin{table}[H]
\centering
\caption{Les deux alertes ajoutées}\label{tab:s6-alertes}
\begin{tabularx}{\linewidth}{@{}P{2.6cm}LL@{}}
\toprule
 & \code{SiteWebIndisponible} & \code{ControlePlanifieEnRetard} \\
\midrule
Condition & Contrôle du site en échec (\code{infraready_controle_succes == 0}) & Dernier contrôle de plus de 300 s, ou métrique absente \\
Durée, gravité & 2 min (deux contrôles consécutifs), critical & 5 min, warning \\
Première action & Journal du contrôle, puis état de nginx & État du timer, puis journal du contrôle \\
Test contrôlé & nginx arrêté de 09:26 à 09:30 & Timer arrêté de 09:54 à 10:04 \\
Résultat & Firing à +3 min 08 ; \code{InstanceDown} reste inactive ; résolue 27 s après & Firing à 10:04:01 ; résolue à 10:04:17 \\
\bottomrule
\end{tabularx}
\end{table}

\illustration{04-etats-alertes.pdf}{0.42}{Cycle de vie des deux alertes pendant les tests}{fig:s6-etats}

\illustration{P-84-grafana-dashboard-site-en-echec.png}{0.78}{Dashboard figé à 09:29:40 : contrôle du site en ÉCHEC et alerte \code{SiteWebIndisponible} active, alors que LINUX01 reste EN LIGNE (\preuve{P-84})}{fig:s6-echec}

\illustration{P-84-prometheus-historique-controle-succes.png}{0.85}{Prometheus : \code{infraready_controle_succes} tombe à 0 pendant l'arrêt de nginx (axe horaire en UTC, soit l'heure de Paris moins 2 h)}{fig:s6-histo}

\illustration{P-83-rendu-timer-journal-metriques.png}{0.95}{Timer, journal du contrôle et métriques publiées sur LINUX01 (rendu de la sortie, \preuve{P-83})}{fig:s6-timer}

% -----------------------------------------------------------------------------
\section{UC-05 — Incident : arrêt brutal du contrôleur de domaine}

\paragraph{Symptôme.} Au démarrage du 23/10, WIN-DC01 journalise les événements 41 (Kernel-Power)
et 6008 (\og l'arrêt précédent du système était inattendu \fg{}). Le même couple d'événements
existe après l'arrêt du 19/10 : l'incident est systématique. \textbf{Impact supposé :} risque
d'incohérence de la base Active Directory et des zones DNS à chaque arrêt de l'hôte.

\illustration{P-80-rendu-win-dc01-evenements-arret.png}{0.95}{Événements 41 et 6008 de WIN-DC01 (rendu de la sortie, \preuve{P-80})}{fig:s6-evt}

\illustration{05-diagnostic-dc.pdf}{0.48}{Démarche de diagnostic de l'incident INC-2026-10-23-01}{fig:s6-diag}

\illustration{06-arret-sequence.pdf}{0.8}{Ce qui se passe à l'arrêt de l'hôte}{fig:s6-arret}

\illustration{P-86-rendu-proxmox-arrets-vm100-reproduction.png}{0.95}{Journal des tâches Proxmox : la VM 100 tuée après 182–183 s, puis la reproduction contrôlée qui échoue sur délai (rendu de la sortie, \preuve{P-86})}{fig:s6-pve}

\begin{itemize}
  \item \textbf{Cause confirmée :} Windows Server ignore la demande d'arrêt ACPI à l'écran d'ouverture
        de session lorsque la stratégie \code{ShutdownWithoutLogon} vaut 0, ce qui est sa valeur par
        défaut ; Proxmox attend 180 s puis tue la VM.
  \item \textbf{Limite mise en évidence :} Prometheus tourne sur OPS01, arrêtée en même temps : la
        supervision ne peut pas observer cet incident.
  \item \textbf{Correction proposée :} passer \code{ShutdownWithoutLogon} à 1 par \code{windows_base.yml}
        (retour arrière : valeur 0), puis valider par un arrêt ACPI réel et l'absence d'événements 41/6008.
\end{itemize}

\attention{Correction non appliquée}{%
La correction modifie le contrôleur de domaine partagé : elle attend une validation explicite de
l'équipe. Rapport complet : \code{docs/rapport-incident.md} (INC-2026-10-23-01).}

% -----------------------------------------------------------------------------
\section{UC-06 — Intégration aux livrables}

\illustration{11-git-ci.pdf}{1}{Du laboratoire au dépôt : branches, demandes de fusion, CI et fusion}{fig:s6-git}

\illustration{P-64-github-pr-seance6-fusionnees.png}{0.8}{Demandes de fusion de la séance 6, fusionnées avec des contrôles CI 8/8 (\preuve{P-64})}{fig:s6-pr}

% -----------------------------------------------------------------------------
\section{UC-C1 et UC-C2 — Maintenance encadrée par un contrôle de santé}

Le playbook \code{controle_sante.yml} contrôle les trois serveurs en lecture seule
(\code{changed=0}, y compris en \code{--check}) et produit un tableau OK / ATTENTION / ECHEC / NON
APPLICABLE / INACCESSIBLE ; il échoue (code 2) s'il reste un échec, ce qui permet de l'utiliser
comme critère d'arrêt d'une maintenance.

\illustration{08-controle-sante.pdf}{1}{Structure du playbook \code{controle_sante.yml}}{fig:s6-cs}

\illustration{P-82-rendu-tableau-controle-sante.png}{0.98}{Contrôle de santé des trois serveurs : 23 OK, 3 ATTENTION, 1 NON APPLICABLE, 0 ECHEC (rendu de la sortie, \preuve{P-80b} / \preuve{P-82})}{fig:s6-cstab}

\illustration{07-maintenance.pdf}{0.5}{Processus de maintenance suivi, avec ses points d'arrêt}{fig:s6-maint}

\illustration{P-87-rendu-maintenance-avant-apres.png}{0.95}{Maintenance de LINUX01 : 3 paquets installés, 5 reportés par le phasage d'Ubuntu, aucun écart de statut avant / après (rendu de la sortie, \preuve{P-87})}{fig:s6-maintres}

% -----------------------------------------------------------------------------
\section{UC-C3 — Observation ponctuelle et suivi dans le temps}

La mesure du jour (OPS01 \code{/} à 26,3 \%) ne dit rien d'une tendance. L'historique de Prometheus
ne couvre que 15 heures sur 91, car le laboratoire était éteint le reste du temps
(figure~\ref{fig:s6-couv}) : OPS01 gagne environ 0,12 Gio sur 15 heures de fonctionnement, LINUX01
reste stable. Une pente calculée par Prometheus sur 45 minutes de données annonçait au contraire
LINUX01 à 4,7 \% dans 30 jours : \textbf{une pente calculée sur si peu de données n'est pas une
tendance}. Une alerte prédictive est proposée, à valider après une semaine de fonctionnement continu.

\illustration{09-couverture-donnees.pdf}{1}{Périodes couvertes par l'historique Prometheus}{fig:s6-couv}

\illustration{P-88-grafana-evolution-stockage-19-23-oct.png}{0.95}{Évolution du stockage du 19 au 23/10 : trois îlots de données ; on y voit aussi \code{C:} de WIN-DC01 passer de 24 \% à 43 \% puis 37 \% pendant Windows Update (\preuve{P-88})}{fig:s6-stock}

% -----------------------------------------------------------------------------
\section{UC-A3 (défi) — Quotas sur le disque de test}

Des quotas ont été préférés à un RAID de démonstration, faute de disques supplémentaires réservés au
TP. Le playbook \code{linux_quotas.yml} vérifie d'abord le numéro de série du disque et ses options
de montage, et s'arrête avant toute modification si ce n'est pas le disque de test.

\illustration{10-quotas.pdf}{0.95}{Comportement des quotas observé}{fig:s6-quotas}

\illustration{P-89-rendu-test-quotas-alice.png}{0.95}{Test réel : alice acceptée à 220 Mio avec délai de grâce, puis bloquée à 250 Mio (rendu de la sortie, \preuve{P-89})}{fig:s6-quotatest}

\textbf{Limites :} root n'est pas soumis aux quotas ; \code{repquota} lit un fichier mis à jour en
différé, alors que \code{quota -u} interroge le noyau. Le premier essai n'était pas idempotent
(\code{repquota} sans \code{-v} masque les comptes sans occupation) : corrigé (\preuve{P-89a} → \preuve{P-89}).
